\chapter{Acompanhamento de Egressos}

O curso de \nomecursonovo da Universidade Federal de Uberlândia reconhece a importância do acompanhamento sistemático de seus egressos como ferramenta para:

\begin{itemize}
    \item Avaliar a efetividade da formação oferecida;
    \item Identificar demandas do mercado de trabalho e tendências tecnológicas;
    \item Promover a atualização do Projeto Pedagógico do Curso (PPC);
    \item Fortalecer o vínculo institucional com os profissionais formados;
    \item Incentivar a participação de ex-alunos(as) em projetos, eventos e processos formativos.
\end{itemize}

\section{Instrumentos e Estratégias de Acompanhamento}

Atualmente, o acompanhamento dos egressos ocorre por meio de:

\begin{itemize}
    \item Aplicação periódica de questionários de egresso, organizados pela Coordenação do Curso, com apoio do Núcleo Docente Estruturante (NDE) do curso e da Comissão Própria de Avaliação (CPA) da UFU, a fim de mapear trajetórias profissionais, empregabilidade, continuidade de estudos e percepção sobre a formação recebida;
    \item Participação de egressos em eventos acadêmicos, como semanas de curso, seminários, bancas de TCC, oficinas de carreira, rodas de conversa e eventos de boas-vindas;
    \item Contato com egressos por meio de redes sociais, grupos de ex-alunos e mailing institucional, com apoio da Direção da FEELT e do setor de comunicação da UFU;
    \item Levantamento de indicadores institucionais (tais como tempo até o primeiro emprego, faixa salarial, área de atuação, envolvimento com pesquisa e inovação) a partir de dados coletados localmente e em parceria com a CPA.
\end{itemize}
    
\section{Objetivos do acompanhamento de egressos}

As informações obtidas por esses mecanismos subsidiam:

\begin{itemize}
    \item A revisão da matriz curricular e o ajuste de conteúdos e metodologias;
    \item A validação do perfil profissional do egresso, conforme as necessidades da sociedade e dos setores produtivos;
    \item O fortalecimento da articulação com o mercado de trabalho e com ex-alunos(as) em posições estratégicas;
    \item O desenvolvimento de políticas institucionais de educação continuada, pós-graduação e formação permanente.
\end{itemize}

A UFU, por meio da CPA e dos colegiados de curso, estimula a construção de políticas institucionais permanentes para o acompanhamento de egressos, integrando bases de dados, experiências bem-sucedidas e iniciativas dos cursos de graduação.